% ============================================================
%  第四部分：策略耦合
% ============================================================

\section{策略耦合：从状态识别到决策姿态}
\label{sec:strategy}

识别出状态只是中间产物。本章回答“识别之后怎么做”：先沿资产配置与 CTA 两条支线梳理学术界验证过的耦合方式，再统一到一个工程问题——\emph{检测延迟与假阳性如何共同决定姿态设计}。

\subsection{资产配置支线：从静态权重到状态条件权重}

\subsubsection{问题的起源：分散化是状态依赖的}

Ang \& Bekaert~\cite{ang2002} 的研究确立了一个后来被反复引用的基本事实：国际股票收益的相关性在高波动（熊市）状态下上升——分散化的收益本身是状态的函数。他们在 regime-switching 框架下求解国际配置问题，展示了状态条件化对配置决策的实质影响。这一支线后续沿着两条路径深化：

\begin{itemize}[itemsep=2pt]
\item \textbf{可行性与稳健性检验}。Bulla et al.~\cite{bulla2011} 系统比较了多种 Markov-switching 配置策略：结论是谨慎的——可盈利策略\emph{可能}存在，但盈利能力高度依赖市场、样本与实施成本，需要逐案验证而非默认可行。
\item \textbf{动态优化化}。Nystrup et al.~\cite{nystrup2018} 用模型预测控制（MPC）替代“固定规则 $+$ 状态开关”：以 HMM 的均值与方差预测为输入，逐期重新优化配置。其含交易成本、1 日执行延迟、零利率现金的样本外结果优于“静态决策规则”，且实施“交易惩罚”（减少交易次数）会进一步提高稳健性。
\end{itemize}

\subsubsection{三种条件化形式及其取舍}

\begin{table}[htbp]
\centering
\footnotesize
\begin{tabular}{>{\raggedright\arraybackslash}p{2.6cm}>{\raggedright\arraybackslash}p{4.7cm}>{\raggedright\arraybackslash}p{4.7cm}}
\toprule
\textbf{形式} & \textbf{机制} & \textbf{适用与代价} \\
\midrule
软路由 & $w_t = \sum_k \pi_t(k)\, w_k$：以状态概率对子权重加权 & 平滑、无跳变；但“概率中间态”的权重组合可能两头不靠 \\
硬切换 $+$ 缓冲 & 概率越过迟滞带（hysteresis）才切换 & 换手低、可审计；需要设计迟滞区间宽度 \\
连续优化（MPC） & 每期以状态预测为输入重新求解优化问题 & 最灵活、天然含约束；实现复杂度最高 \\
\bottomrule
\end{tabular}
\caption{状态条件化配置的三种实现形式}
\label{tab:allocation-forms}
\end{table}

\subsubsection{防御状态下的组合理论依据}

第~\ref{sec:frontier}~节引用的传染分析~\cite{scherer2020} 在此给出配置含义：当“高波动防御”状态伴随策略间传染时，分散化收益消失，杠杆成为主要风险暴露——最优回应是\emph{系统性降低总杠杆}并回避对传染敏感的资产。这为“防御状态下机械降杠杆”提供了理论支撑：它不是胆小，而是约束下最优解的结构性变化。

\subsection{CTA 支线：趋势跟踪与状态}

\subsubsection{证据基础}

趋势跟踪是跨资产（期货）策略中证据最强的家族：Moskowitz et al.~\cite{moskowitz2012} 在 58 个流动性市场中记录了 1--12 月的时序动量；Hurst et al.~\cite{hursto2017} 把证据延伸到 1880--2016 的百年样本，显示各年代风险调整收益均为正且与股债相关性极低；其机构版拆解~\cite{hurst2013} 表明 CTA 基金收益的大部可由简单时序动量策略解释。在信号层面，Levine \& Pedersen~\cite{levine2016} 对不同趋势测度的横向对比表明：“用什么趋势”与“是否用趋势”同等重要。

\subsubsection{状态依赖性是趋势跟踪研究的新重心}

既然趋势家族长期有效，为什么还需要 Regime？三个近年的答案：

\begin{itemize}[itemsep=2pt]
\item \textbf{转折点密度决定短期成败}。Garg et al.~\cite{garg2024} 量化了“转折点（快慢动量符号翻转）密度”与趋势策略绩效的负相关：转折频繁时期，标准趋势规则系统性亏损。这解释了趋势策略的“难受期”不是随机分布，而是聚集在震荡状态。
\item \textbf{最优信号速度是状态相关的}。Goulding et al.~\cite{goulding2023} 的“动量转折点”框架显示，转折点后应提高快信号权重——即最优动量参数随状态切换。
\item \textbf{更根本的：给定状态，标准规则并非最优}。Zakamulin~\cite{zakamulin2026} 把“信号”与“配置”分离：在状态已被识别的条件下，从第一性原理导出每个状态的 Sharpe 最优暴露。标准时序动量规则（涨多、跌空）被证明是这一框架的\emph{特例}，且一般是\emph{次优}的——一个被反复验证的反直觉结论：\textbf{熊市全额做空往往远离最优，最优暴露常接近零（甚至轻微为正）}。跨 18 个组合数据集的样本外检验中，regime-optimal 配置全面优于标准时序动量与动态速度动量（两状态平均 Sharpe 从 $0.208$ 升至 $0.506$；四状态从 $0.496$ 升至 $0.628$）。
\end{itemize}

\begin{remark}[结论的含义]
它的含义超出了参数调优：趋势跟踪的最优形式可能不是“更聪明的信号”，而是“在已能识别的状态里更聪明地配置”。这与第~\ref{sec:frontier}~节的综合判断完全同构——\emph{系统的进步空间主要位于状态条件化与姿态设计，而非识别精度的军备竞赛}。
\end{remark}

\subsubsection{因子层面的条件化}

状态切换不止改变仓位，也改变因子有效性：Wang et al.~\cite{wang2020} 用 HMM 识别美股状态并在不同状态下轮换因子模型（价值/质量/动量等），样本外优于任何单一因子模型。这条路线与企业级实践（第~\ref{sec:engineering}~节体系中对下游因子系统的约束输出：偏好因子类别随状态切换、ML 因子权重上限随状态下调）一致：\textbf{状态是因子权重方案的元参数}。

\subsubsection{波动率目标化的状态化}

波动率目标化本身已经是“状态感知”的（以其方式）：仓位与已实现波动率成反比。把它与 Regime 结合的自然形式是：不同状态下使用不同的目标波动率与调整速度——防御状态下目标波动率下移、调整更快（快速降杠杆）；趋势友好状态下目标波动率上调、调整更慢（避免被短期噪声震出）。

\subsection{延迟与假阳性的姿态设计：把下界变成设计参数}

第一部分证明的“延迟—误报前沿”在此转化为具体设计问题。两类错误的成本是不对称的：

\begin{itemize}[itemsep=2pt]
\item \textbf{假阴性（切换后确认太晚）}：错过新状态初期的机会/风险，代价 = 延迟 $\times$ 单期错误暴露的损失；
\item \textbf{假阳性（误报切换）}：无谓的调仓磨损 $+$ 在旧状态中的错误姿态，代价 = 换手成本 $+$ 错误期数的损失。
\end{itemize}

工程上由此导出四条设计原则：

\begin{enumerate}[itemsep=3pt]
\item \textbf{非对称响应}：进入防御可以快（假阴性在防御方向上的代价更高——回撤不可逆），恢复进攻必须慢（假阳性在进攻方向上的代价更高——反复假突破磨损）。这一不对称是多数高质量系统的隐含设计。
\item \textbf{连续映射优先于二值开关}：用状态概率的连续函数映射仓位（软路由），在中间态自然降速，减少高频抖动。
\item \textbf{确认机制显式化}：投票（多渠道证据）、确认期（连续 $N$ 期越过阈值）、迟滞带——这些都是把“工作点”写成代码的方式。第~\ref{sec:engineering}~节体系的“三窗口 HMM 投票 $+$ 规则法交叉”正是这一原则的实现。
\item \textbf{把延迟写进回测}：任何状态条件化策略的回测必须包含 1--2 期的信号—执行延迟（参照 Nystrup~\cite{nystrup2018} 与 Shu~\cite{shu2024a} 的实验设计），否则延迟成本被系统性低估。
\end{enumerate}

\begin{table}[htbp]
\centering
\footnotesize
\begin{tabular}{p{3.0cm}p{4.2cm}p{4.8cm}}
\toprule
\textbf{设计维度} & \textbf{保守设置} & \textbf{适用条件} \\
\midrule
切换阈值 & 高（需 2/3 证据一致） & 换手成本高、状态持续期长 \\
确认期 & 连续 $N$ 期越过 & 单期噪声大（高频数据） \\
迟滞带宽度 & 宽（进入/退出阈值分开） & 状态边界模糊（震荡 $\leftrightarrow$ 弱趋势） \\
响应不对称 & 防御快、进攻慢 & 回撤约束严格（含杠杆组合） \\
\bottomrule
\end{tabular}
\caption{延迟—误报工作点的工程化旋钮}
\label{tab:workpoint}
\end{table}

\keynote{识别之后的一切，可以归结为把“状态概率”翻译成“仓位姿态”。优秀系统与平庸系统的差别不在翻译公式的复杂度，而在于对两类错误的非对称处理：进入防御可以快，恢复进攻必须慢。}